\documentclass[11pt,letterpaper]{article}
\usepackage[T1]{fontenc}
\usepackage{lmodern,amsmath,amssymb,amsthm,mathtools,mathrsfs}
\usepackage[margin=1in,headheight=15pt]{geometry}
\usepackage{microtype,booktabs,longtable,array,xcolor,fancyhdr,enumitem}
\usepackage[colorlinks=true,linkcolor=blue!40!black,citecolor=blue!40!black,urlcolor=blue!40!black]{hyperref}
\usepackage{bookmark}
\allowdisplaybreaks[2]
\emergencystretch=2em
\numberwithin{equation}{section}
\newtheorem{theorem}{Theorem}[section]
\newtheorem{lemma}[theorem]{Lemma}
\newtheorem{proposition}[theorem]{Proposition}
\newtheorem{corollary}[theorem]{Corollary}
\theoremstyle{definition}\newtheorem{definition}[theorem]{Definition}
\theoremstyle{remark}\newtheorem{remark}[theorem]{Remark}
\DeclareMathOperator{\Li}{Li}
\DeclareMathOperator{\CT}{CT}
\DeclareMathOperator{\FP}{FP}
\DeclareMathOperator{\wt}{wt}
\newcommand{\Q}{\mathbb Q}
\newcommand{\N}{\mathbb N}
\newcommand{\mz}[1]{\zeta(#1)}
\newcommand{\sref}[1]{\texttt{#1}}
\newcommand{\dd}{\,d}
\newcommand{\sha}{\mathbin{\mathrm{sh}}}
\newcommand{\shuffle}{\mathrm{sh}}
\pagestyle{fancy}\fancyhf{}
\fancyhead[L]{\small Mixed tails, zeta anchors, and Stieltjes transport}
\fancyhead[R]{\small ProveIt continuation}
\fancyfoot[C]{\thepage}
\setlist[enumerate]{leftmargin=*,itemsep=5pt}
\hypersetup{pdftitle={Mixed Hurwitz-Tail Moments, Zeta-Anchor Reductions, and Stieltjes Transport},pdfauthor={Research continuation prepared for Vladimir Reshetnikov},pdfsubject={Exact polylogarithm, multiple-zeta, polygamma, Gamma and Stieltjes identities}}
\begin{document}
\begin{titlepage}
\vspace*{10mm}
{\sffamily\large PROVEIT RESEARCH CONTINUATION\par}
\vspace{12mm}
{\LARGE\bfseries Mixed Hurwitz-Tail Moments,\par
Zeta-Anchor Reductions,\par
and Stieltjes Transport\par}
\vspace{8mm}
{\large Exact identities, a parity-refined finite compiler,\par
and explicit rational certificates\par}
\vspace{12mm}
{Prepared for Vladimir Reshetnikov\par}
\smallskip
{11 October 2026\par}
\vspace{10mm}
\begin{abstract}
We continue the harmonic-tail and Gamma--Stieltjes branches of the ProveIt
polylogarithm programme. Three explicitly defined anchors in the existing
fifth- and sixth-power moment generators are reduced to ordinary zeta
products and, at weight eight, a single specified depth-two value.
The reductions are proved by finite rational combinations of classical
double-shuffle relations; complete short certificates are printed and
independently replayed.

For an arbitrary product of integer-order Hurwitz tails, including odd
and unequal orders, we give a terminating formula for every regular
centered moment. The answer lies in the rational span of admissible
multiple zeta values of weight at most $W-1$, where $W$ is the sum of
the tail orders. For odd centered degree the ceiling improves to $W-2$.
A separate, simpler formula treats the full $(2,3)$ family and gives
an ordinary polynomially subtracted series with value exactly $-15$;
we supply a direct finite telescoper for it.

Finally, weighting by successive differences of a tail product gives
all-index generalized-Stieltjes sums, polygamma identities, anchored
Gamma-ratio sums, and primitives in terms of Hurwitz order derivatives
at zero. A finite multiple-polylogarithm kernel identifies their
coefficients with convergent final-index multiple-zeta derivatives.
The article distinguishes these continuations from the already proved
all-power even-tail recurrences. It includes source-status notes,
an independently replayable exact certificate suite, numerical
diagnostics, and a precise further-research agenda.
\end{abstract}
\vfill
\noindent\small\textbf{Scope.} The claims are relative to the explicitly
inspected source files, not an exhaustive priority claim over the
literature or the repository. No arithmetic independence, Gaussian
$S_6/S_8$ identity, or proof-assistant formalization is asserted.
\end{titlepage}
\tableofcontents
\clearpage

\section{Results, sources, and conventions}
\label{sec:scope}

The starting point is a collection of proved identities, not a blank
research programme. The quartic reconstruction report \cite{Quartic}
already evaluates all regular centered quartic trigamma moments.
The later Bernoulli-cancellation report \cite{Bernoulli} proves the
all-power recurrence, gives fifth- and sixth-power generators, and
extends the mechanism to repeated even-order tails. The lattice-contact
report \cite{Lattice} further treats arbitrary unequal even-order
products. Those theorems must not be presented as new here.

The present continuation makes three different contributions. First,
it reduces the named higher anchors $J_7,K_8,K_9$ explicitly. Second,
it gives a cutoff-normalized compiler for arbitrary integer tail orders,
with a parity-refined weight ceiling and a particularly simple odd-order
mixed family. Third, it connects the successive differences of these
tail products to Stieltjes constants, Gamma ratios, and order derivatives
of multiple zeta functions. The transport identities use classical
summation by parts and shift recurrences; the contribution is their
explicit combined realization and normalization, not the underlying
summation principle.

The most elementary new specialization is worth displaying first.
Put
\begin{equation}\label{eq:AB}
 A_n=\zeta(2,n+1)=\psi'(n+1),\qquad
 B_n=\zeta(3,n+1)=-\frac12\psi^{(2)}(n+1),\qquad
 x_n=n+\frac12.
\end{equation}
Then
\begin{equation}\label{eq:intro-rational}
\boxed{\displaystyle
 \sum_{n=0}^{\infty}\left[
 (23x_n+232x_n^3+240x_n^5)A_nB_n-120x_n^2-76
 \right]=-15.}
\end{equation}
The bracket is a single absolutely summable quantity. None of its
polynomial pieces is summed separately. Section~\ref{sec:pair} proves
this by an explicit rational telescoper, independently of the
multiple-zeta reduction software.

\subsection{A scoped source audit}

The core anchor definitions were read at the immutable repository ref
\sref{0dfef47533f8}; the supplementary unequal-even-tail report was
inspected at \sref{cc569dca8622}. Full immutable identifiers are
recorded in the provenance file and the pinned bibliography links.
The earlier quartic source and the canonical status files were also
read. Exact paths and the scope of each read are in
\sref{SOURCE\_AUDIT.md} and \sref{provenance.json}.

The incoming directory listing and its intake instructions were
inspected, but its binary archives were not exhaustively extracted or
audited in this session. In particular, absence from the inspected
passages is not evidence of absence from every incoming report.
The article therefore offers explicit proved continuations and
integration candidates, not a claim that every surrounding open
question has been surveyed or resolved.

Classical ingredients are credited: Hurwitz asymptotics and Gamma
relations \cite{DLMF}, products of zeta tails ordered by their minimum
index \cite{Hoffman}, and double-shuffle relations \cite{IKZ}.
The existing fifth- and sixth-power generators and their original
formal-coordinate ranks retain their source attribution.

\subsection{Notation}

We use descending-index multiple zeta values
\[
 \zeta(s_1,\ldots,s_d)
 =\sum_{n_1>\cdots>n_d\ge1}
     \frac1{n_1^{s_1}\cdots n_d^{s_d}}.
\]
An integer word is \emph{admissible} when $s_1\ge2$ and all other
indices are positive. Its weight is $s_1+\cdots+s_d$. The empty word
has zeta value $1$. Products can always be expanded into a finite
linear combination of MZVs by the harmonic, or stuffle, product.
No independence of the resulting numerical values is assumed.

Bernoulli polynomials satisfy
\[
 \frac{te^{xt}}{e^t-1}=\sum_{j\ge0}B_j(x)\frac{t^j}{j!},
 \qquad B_j=B_j(0),\qquad B_1=-\frac12.
\]
For generalized Stieltjes constants our convention is
\begin{equation}\label{eq:stieltjes-convention}
 \zeta(1+u,a)=\frac1u+
 \sum_{q\ge0}\frac{(-1)^q}{q!}\gamma_q(a)u^q,
 \qquad \gamma_0(a)=-\psi(a).
\end{equation}
Order derivatives of zeta functions are always distinguished from
argument derivatives. Positive real shifts suffice for every
Stieltjes and Gamma identity below. This avoids any implicit
choice of logarithmic branches.

\section{Explicit reduction of the higher moment anchors}
\label{sec:anchors}

For a composition $(a_1,\ldots,a_\ell)$ of $e$, write
$(a_1,\ldots,a_\ell)\vDash e$. Define the convergent anchor
\begin{equation}\label{eq:anchor}
 \mathcal Z_{e,h}=
 \sum_{(a_1,\ldots,a_\ell)\vDash e}
 \frac{e!}{a_1!\cdots a_\ell!}
 \zeta(2a_1,\ldots,2a_{\ell-1},2a_\ell-h).
\end{equation}
Only the negative $h$ instances occurring below are needed here.
Ordering the equal and strictly ordered tail indices gives the
ordinary identity
\begin{equation}\label{eq:anchor-drop}
 \mathcal Z_{e,h}
 =\sum_{n=1}^{\infty}n^h(A_{n-1}^e-A_n^e),\qquad h\le0.
\end{equation}
This ordering argument is classical \cite{Hoffman}; the explicit
reductions that follow are the additional step.

The Bernoulli-cancellation source defines
\begin{align}
 J_7&=\mathcal Z_{2,-3}=2\zeta(2,5)+\zeta(7),\label{eq:Jdef}\\
 K_8&=\frac56\mathcal Z_{3,-2}+\frac{289}{32}\zeta(8),\label{eq:K8def}\\
 K_9&=\mathcal Z_{2,-5}-5\mathcal Z_{4,-1}.\label{eq:K9def}
\end{align}
In particular,
\[
 \mathcal Z_{3,-2}
 =6\zeta(2,2,4)+3\zeta(2,6)+3\zeta(4,4)+\zeta(8).
\]
These definitions are inherited, rather than newly introduced
unknown constants.

\begin{theorem}[Three explicit anchor reductions]\label{thm:anchors}
The anchors in \eqref{eq:Jdef}--\eqref{eq:K9def} satisfy
\begin{align}
 J_7={}&21\zeta(7)-8\zeta(2)\zeta(5)
                                  -4\zeta(3)\zeta(4),\label{eq:J7}\\
 K_8={}&10\zeta(2)\zeta(3)^2-50\zeta(3)\zeta(5)
                -10\zeta(5,3)+\frac{4927}{96}\zeta(8),\label{eq:K8}\\
 K_9={}&1086\zeta(9)-222\zeta(2)\zeta(7)
                +171\zeta(3)\zeta(6)-858\zeta(4)\zeta(5).
 \label{eq:K9}
\end{align}
These are identities of convergent real series. They do not assert
that $\zeta(5,3)$ is irreducible in any larger class of constants.
\end{theorem}

\begin{proof}
Expand the left sides by \eqref{eq:anchor}, and expand products on
the right by the stuffle product. At each weight the difference
is a rational vector on the finite set of admissible words.
Appendix~\ref{app:certificates} specifies a finite family of true
double-shuffle and duality relations and gives every coefficient
expressing these three difference vectors as sums of those relations.
The certificates use respectively $3$, $17$, and $54$ nonzero rows.
Their rational vector identities are checked independently by
\sref{verification/exact.py}; the complete row convention and
coefficient tables are also printed in the appendix.

Each row is a genuine MZV identity, and the stated rational linear
combination is exactly the target difference, so the difference
vanishes. This proof uses neither integer-relation recognition nor
a conjectural completeness theorem for double shuffle.
\end{proof}

For comparison, the exact computation gives the convenient equivalent
anchor values
\begin{align}
 \mathcal Z_{3,-2}
   &={}12\zeta(2)\zeta(3)^2-60\zeta(3)\zeta(5)
           -12\zeta(5,3)+\frac{203}{4}\zeta(8),\label{eq:Z3}\\
 \mathcal Z_{4,-1}
   &={}42\zeta(2)\zeta(7)-35\zeta(3)\zeta(6)
           +170\zeta(4)\zeta(5)-210\zeta(9).\label{eq:Z4}
\end{align}
The first follows immediately from \eqref{eq:K8def} and
\eqref{eq:K8}. For the second, the classical odd-weight double-zeta
reduction, or the same finite row system, gives
\[
 \mathcal Z_{2,-5}
 =36\zeta(9)-12\zeta(2)\zeta(7)-4\zeta(3)\zeta(6)
                                    -8\zeta(4)\zeta(5).
\]
Insert this into \eqref{eq:K9def}. The auxiliary reduction can also
be reconstructed with the included exact reducer; it is not inferred
from a decimal fit.

\subsection{Consequences for the existing quintic and sextic generators}

Retain the inherited elementary combinations
\begin{align}
 Q_5&=4\zeta(2)\zeta(3)-10\zeta(5),\label{eq:Q5}\\
 Q_7&=51\zeta(3)\zeta(4)+10\zeta(2)\zeta(5)
                              -\frac{301}{4}\zeta(7).\label{eq:Q7}
\end{align}
The source proves $\sum_{n\ge0}A_n^4=Q_7$ and
\begin{equation}\label{eq:inherited-M51}
 \sum_{n\ge0}x_nA_n^5=-K_8+\frac52Q_7.
\end{equation}
Thus the additional arithmetic reduction yields
\begin{equation}\label{eq:quintic-standard}
\begin{split}
 \sum_{n\ge0}x_nA_n^5={}&
 \frac{255}{2}\zeta(3)\zeta(4)+25\zeta(2)\zeta(5)
                       -\frac{1505}{8}\zeta(7)\\
 &-10\zeta(2)\zeta(3)^2+50\zeta(3)\zeta(5)
                       +10\zeta(5,3)-\frac{4927}{96}\zeta(8).
\end{split}
\end{equation}
The inherited moment relation \eqref{eq:inherited-M51} is explicitly
credited; its expansion into the standard coordinates in
\eqref{eq:quintic-standard} is the continuation.

For clarity, let $I^q$ mean the $q$-fold primitive anchored at zero,
and let $H_d=2\sinh(t/2)G_d$ be the source's regular moment generator.
Its sextic formula is
\begin{equation}\label{eq:inherited-H6}
\begin{split}
 H_6(t)={}&R_{11}t+6I^2(\cosh(t/2)G_5(t))
                 -\frac52I^4H_4(t)+\frac12I^8H_2(t)\\
 &+\frac{K_9}{12}t^3+\frac{J_7}{240}t^5-\frac{Q_5}{10080}t^7,
 \qquad R_{11}=\mathcal Z_{6,1}.
\end{split}
\end{equation}
Equations \eqref{eq:J7} and \eqref{eq:K9} eliminate the two named
lower anchors here. The weight-eleven anchor $R_{11}$ is not reduced
in this article. Nor does rewriting the coordinates prove independence
of the source's eleven numerical values.

\section{Arbitrary integer-order tail products and regular parity}
\label{sec:mixed}

Fix a nonempty tuple $\boldsymbol p=(p_1,\ldots,p_d)$ of integers
$p_i\ge2$. Define
\begin{equation}\label{eq:mixed-def}
 T_n=\prod_{i=1}^d\zeta(p_i,n+1),\qquad
 W=\sum_{i=1}^dp_i,\qquad R=W-d,\qquad x_n=n+\frac12.
\end{equation}
Here $R\ge1$. Odd orders are allowed. For integer $p\ge2$,
\[
 \zeta(p,n+1)=\frac{(-1)^p}{(p-1)!}\psi^{(p-1)}(n+1),
\]
so every result is simultaneously an identity for products of
polygamma functions.

Set
\begin{equation}\label{eq:etas}
 \eta_{p,0}=\frac1{p-1},\qquad
 \eta_{p,j}=\frac{(p)_{2j-1}}{(2j)!}B_{2j}(1/2)\quad(j\ge1),
 \qquad
 c_k=\sum_{k_1+\cdots+k_d=k}\prod_i\eta_{p_i,k_i}.
\end{equation}
The centered Hurwitz expansion \cite{DLMF} is
\begin{equation}\label{eq:centered-expansion}
 T_n\sim\sum_{k\ge0}c_kx_n^{-R-2k}.
\end{equation}
It is an asymptotic expansion, not a convergent series at a fixed $n$.

\begin{proposition}[Poles and ordinary regular moments]\label{prop:regular}
The series
\[
 \mathcal D_{\boldsymbol p}(s)=\sum_{n\ge0}T_nx_n^{-s}
\]
converges absolutely for $\Re s>1-R$ and continues meromorphically
to the plane. Its poles are exactly the simple poles
$s=1-R-2k$, with residue $c_k\ne0$.
For each nonnegative integer $m\equiv R\pmod2$, the regular value
$M_{\boldsymbol p,m}=\mathcal D_{\boldsymbol p}(-m)$ is the ordinary
absolutely convergent sum
\begin{equation}\label{eq:regular-sum}
 \boxed{M_{\boldsymbol p,m}=
 \sum_{n=0}^{\infty}\left[
 x_n^mT_n-
 \sum_{0\le k\le(m-R)/2}c_kx_n^{m-R-2k}\right].}
\end{equation}
A sum whose upper bound is negative is empty.
\end{proposition}

\begin{proof}
A finite truncation of \eqref{eq:centered-expansion}, with its
remainder estimate, writes $\mathcal D$ as a normally convergent
remainder series plus finitely many functions
$c_k\zeta(s+R+2k,1/2)$. Increasing the truncation continues it to
successively larger left half-planes. The pole and residue assertions
follow. Since $\eta_{p,j}$ has sign $(-1)^j$, every term in the
convolution for $c_k$ has sign $(-1)^k$; hence $c_k$ is nonzero.

At $s=-m$ in the indicated parity, every removed term has added
Hurwitz value $\zeta(-2j,1/2)=0$. The remainder in the bracket
is $O(x_n^{-2})$, or faster when no subtraction is needed. This
proves \eqref{eq:regular-sum}, with no unspecified summation of
divergent polynomial series.
\end{proof}

\begin{theorem}[Finite mixed-order reduction and refined weight ceiling]
\label{thm:compiler}
For every tuple in \eqref{eq:mixed-def} and every regular degree $m$,
there is a terminating rational algorithm expressing
$M_{\boldsymbol p,m}$ as a linear combination of $1$ and admissible
MZVs of weight at most $W-1$. If $m$ is odd, the bound improves to
$W-2$. The coefficient algorithm in Section~\ref{sec:compiler}
specifies the value without choosing a conjectural basis of MZVs.
\end{theorem}

The all-even case overlaps the source's Bernoulli recurrences, which
already give finite-span conclusions. The theorem's useful additional
features are its allowance of arbitrary odd orders, its explicit
cutoff-normalized finite procedure, and the uniform odd-degree
refinement. The weight bound is a spanning statement, not a claim
that a given value has that minimal weight or depth.

\section{A finite-cutoff proof and an executable compiler}
\label{sec:compiler}

\subsection{Harmonic words and their asymptotic coefficients}

For a word $\boldsymbol w=(w_1,\ldots,w_\ell)$ put
\[
 h_{\boldsymbol w}(N)=
 \sum_{N>n_1>\cdots>n_\ell\ge1}
       n_1^{-w_1}\cdots n_\ell^{-w_\ell},\qquad
 h_{\varnothing}(N)=1.
\]
This is an \emph{exclusive} cutoff. Also write
$H_n(\boldsymbol w)=h_{\boldsymbol w}(n+1)$.
First suppose every $w_i\ge2$. Then
\begin{equation}\label{eq:h-asymptotic}
 h_{\boldsymbol w}(N)\sim\sum_{h\ge0}a_{\boldsymbol w,h}N^{-h},
 \qquad a_{\boldsymbol w,0}=\zeta(\boldsymbol w).
\end{equation}
There are no logarithmic terms in this expansion.

Let $z_{q,h}$ denote the coefficient of $N^{-h}$ in
$\zeta(q,N)$ for $q\ge2$. Explicitly,
\begin{equation}\label{eq:hurwitz-coeff}
 z_{q,h}=
 \begin{cases}
 (q-1)^{-1},&h=q-1,\\
 1/2,&h=q,\\
 B_{2j}(q)_{2j-1}/(2j)!,&h=q+2j-1,\ j\ge1,\\
 0,&\text{otherwise}.
 \end{cases}
\end{equation}
The initial conditions and recurrence
\begin{equation}\label{eq:a-recursion}
 a_{\varnothing,0}=1,\quad a_{\varnothing,h}=0\ (h>0),\qquad
 a_{(s,\boldsymbol v),h}
 =-\sum_{r\ge0}a_{\boldsymbol v,r}z_{s+r,h}\quad(h>0)
\end{equation}
compute every required coefficient by finite arithmetic.
For fixed $h$ only finitely many $r$ contribute.

\begin{lemma}\label{lem:a}
Equations \eqref{eq:h-asymptotic}--\eqref{eq:a-recursion} are valid
as asymptotic expansions with remainders to every fixed order.
For $h>0$, $a_{\boldsymbol w,h}$ is a rational linear combination
of zeta values of proper suffixes of $\boldsymbol w$.
\end{lemma}
\begin{proof}
For a nonempty word,
\[
 \zeta(s,\boldsymbol v)-h_{(s,\boldsymbol v)}(N)
 =\sum_{n\ge N}n^{-s}h_{\boldsymbol v}(n).
\]
Insert any finite asymptotic expansion for the shorter word.
Since $s\ge2$, its remainder can be summed absolutely and gives
an error of the next required order. The retained terms are
$a_{\boldsymbol v,r}\zeta(s+r,N)$. Expanding each of them gives
\eqref{eq:a-recursion}. Induction on the word length proves all
assertions, including the suffix statement.
\end{proof}

\subsection{Cutoff constants with a nonpositive or unit first index}

The finite summation step will introduce a possibly nonpositive first
index, but all following indices remain at least two. Let
$\mathcal C(\boldsymbol w)$ denote the coefficient of
$N^0(\log N)^0$ in the expansion of $h_{\boldsymbol w}(N)$.
For an admissible word it is simply $\zeta(\boldsymbol w)$.
The following finite formulas define all other cases that are needed.

For a suffix $\boldsymbol v=(v_1,\ldots,v_\ell)$ with $v_i\ge2$,
\begin{equation}\label{eq:C-one}
\begin{split}
 \mathcal C(1,\boldsymbol v)={}&\gamma\zeta(\boldsymbol v)
 -\sum_{j=1}^{\ell}\zeta(v_1,\ldots,v_j,1,v_{j+1},\ldots,v_\ell)\\
 &-\sum_{j=1}^{\ell}\zeta(v_1,\ldots,v_j+1,\ldots,v_\ell).
\end{split}
\end{equation}
For the empty suffix this means $\mathcal C(1)=\gamma$.
All the other displayed MZVs converge.

For $r\ge0$ define the Faulhaber polynomial
\begin{equation}\label{eq:Faulhaber}
 P_r(X)=\sum_{k=1}^{X}k^r
 =\frac{B_{r+1}(X+1)-B_{r+1}(1)}{r+1}
 =\sum_{\ell\ge1}p_{r,\ell}X^\ell.
\end{equation}
The expression is polynomial notation, not a restriction that $X$
be an integer. If the suffix is empty, set
\begin{equation}\label{eq:C-empty}
 \mathcal C(-r)=P_r(-1)=-\mathbf1_{r=0}.
\end{equation}
If the suffix is $(s,\boldsymbol u)$, set
\begin{equation}\label{eq:C-negative}
 \mathcal C(-r,s,\boldsymbol u)
 =\CT_N\!\left[P_r(N-1)h_{(s,\boldsymbol u)}(N)\right]
   -\sum_{\ell\ge1}p_{r,\ell}
             \mathcal C(s-\ell,\boldsymbol u).
\end{equation}
The constant term of a polynomial times $h$ is obtained directly
from the coefficients in \eqref{eq:a-recursion}. Each recursive
call in \eqref{eq:C-negative} has shorter length.

\begin{lemma}[Exact meaning of the cutoff constants]\label{lem:C}
Equations \eqref{eq:C-one}--\eqref{eq:C-negative} equal the actual
asymptotic cutoff constants. Their recursion terminates and introduces
only rational numbers, admissible MZVs, and a term linear in $\gamma$.
\end{lemma}
\begin{proof}
The finite stuffle identity for $h_1(N)h_{\boldsymbol v}(N)$
isolates $h_{1,\boldsymbol v}(N)$. All other terms converge;
$h_1(N)=\log N+\gamma+O(N^{-1})$ and
$h_{\boldsymbol v}(N)=\zeta(\boldsymbol v)+O(N^{-1})$.
Taking the constant term gives \eqref{eq:C-one}.

For a nonpositive first index, reverse the finite sums:
\[
 h_{(-r,s,\boldsymbol u)}(N)
 =P_r(N-1)h_{(s,\boldsymbol u)}(N)
       -\sum_{\ell\ge1}p_{r,\ell}h_{(s-\ell,\boldsymbol u)}(N).
\]
This is an exact identity for every positive integer $N$.
Taking cutoff constants gives \eqref{eq:C-negative}. The empty
suffix gives $P_r(N-1)$ itself, whose constant is $P_r(-1)$.
Induction proves termination and the stated form of the answer.
\end{proof}

\subsection{The master formula}

Multiplying the finite identities
$\zeta(p,n+1)=\zeta(p)-H_n^{(p)}$ gives
\begin{equation}\label{eq:tail-harmonic-expansion}
 T_n=\sum_{\boldsymbol w}b_{\boldsymbol w}H_n(\boldsymbol w).
\end{equation}
All entries in every nonempty $\boldsymbol w$ are at least two.
The coefficients $b_{\boldsymbol w}$ are finite polynomials in
ordinary zeta values. They are generated by the stuffle product,
starting from $b_{\varnothing}=1$ and multiplying one factor
$\zeta(p)-H_n^{(p)}$ at a time. Every homogeneous term satisfies
$\wt(b_{\boldsymbol w})+\wt(\boldsymbol w)=W$.

Define the centered discrete primitive
\begin{equation}\label{eq:Fm}
 F_m(X)=\frac{B_{m+1}(X+1/2)-B_{m+1}(1/2)}{m+1}
       =\sum_{h\ge1}f_{m,h}X^h.
\end{equation}
It has zero constant term and
$F_m(n+1)-F_m(n)=x_n^m$.

\begin{theorem}[Finite coefficient formula]\label{thm:master}
For the regular degrees in Proposition~\ref{prop:regular},
\begin{equation}\label{eq:master}
 \boxed{M_{\boldsymbol p,m}=
 \sum_{\boldsymbol w\ne\varnothing}b_{\boldsymbol w}
 \sum_{h\ge1}f_{m,h}
 \left[a_{\boldsymbol w,h}
 -\mathcal C(w_1-h,w_2,\ldots,w_\ell)\right].}
\end{equation}
All sums are finite. After the stuffle product is expanded, the
coefficient of $\gamma$ is zero.
\end{theorem}

\begin{proof}
For a nonempty word, reversing the finite sums gives
\begin{equation}\label{eq:finite-centered}
 \sum_{n=0}^{N-1}x_n^m H_n(\boldsymbol w)
 =F_m(N)h_{\boldsymbol w}(N)
 -\sum_{h\ge1}f_{m,h}h_{(w_1-h,w_2,\ldots,w_\ell)}(N).
\end{equation}
Indeed, the contribution with largest harmonic index $k$ has the
inner sum $\sum_{n=k}^{N-1}x_n^m=F_m(N)-F_m(k)$.
Insert \eqref{eq:tail-harmonic-expansion} and take cutoff constants.
The empty word contributes $F_m(N)$, whose constant is zero.
Lemmas~\ref{lem:a} and \ref{lem:C} give precisely \eqref{eq:master}.

It remains to identify this constant with the ordinary sum.
Every subtracted monomial in \eqref{eq:regular-sum} is an even power
of $x_n$. Its finite sum is $F_{2j}(N)$, whose constant is zero.
The sum of the remaining bracket tends to the ordinary value, with
an $O(N^{-1})$ remainder. Thus the original cutoff constant is that
value.

Finally, regular parity makes every exponent in the centered
asymptotic summand even. Its partial sums have no logarithmic term.
In the finite reduction, every occurrence of $\gamma$ comes from
$h_1(N)=\log N+\gamma+o(1)$ and carries exactly the same coefficient
as the corresponding logarithm. That total coefficient is zero.
This proves the cancellation without treating $\gamma$ as an
independent transcendental number.
\end{proof}

\begin{proof}[Proof of Theorem~\ref{thm:compiler}]
The recursions decrease either the suffix length or a bounded
coefficient index, so they terminate. For $h>0$, the term
$a_{\boldsymbol w,h}$ contains only proper suffixes and hence
has weight at most $\wt(\boldsymbol w)-2$.
If $w_1-h\ge1$, the non-$\gamma$ terms in the corresponding
cutoff constant have weight $\wt(\boldsymbol w)-h$.
If $w_1-h\le0$, the negative-index recursion gives weights at
most the weight of the remaining suffix, again no more than
$\wt(\boldsymbol w)-2$. This last bound follows inductively
from \eqref{eq:C-negative}; its polynomial boundary term uses
only suffix values, and each subsequent positive leading index
is reduced rather than increased.

Multiplying by $b_{\boldsymbol w}$ and expanding products preserves
total weight. Since $h\ge1$, this gives the bound $W-1$.
When $m$ is odd, $F_m$ is an even polynomial with zero constant
term, so every nonzero $f_{m,h}$ has $h\ge2$. The same argument
then gives $W-2$.
\end{proof}

\section{The complete odd-degree (2,3) family}
\label{sec:pair}

The general compiler simplifies sharply for the product in
\eqref{eq:AB}. Its total weight is $5$, so all its regular moments
lie in $\operatorname{span}_{\Q}\{1,\zeta(2),\zeta(3)\}$.
A direct formula computes these three coefficients without a
general MZV reduction.

Let $P(X)=\sum_{j\ge1}p_{2j}X^{2j}$ be any even rational polynomial
with $P(0)=0$, and set
\[
 W_P(x)=P(x+1/2)-P(x-1/2),\qquad
 \mathscr M(P)=\sum_{m\text{ odd}}[x^m]W_P(x)\,M_{(2,3),m}.
\]
The last expression means one ordinary bracketed sum with the
subtractions prescribed in \eqref{eq:regular-sum}.
Put
\begin{equation}\label{eq:pair-L}
 \lambda(P)=\sum_{j\ge1}p_{2j}B_{2j-2},\qquad
 Q(X)=P(X)-p_2X^2,
\end{equation}
and define rational polynomials by finite power sums:
\begin{equation}\label{eq:pair-RS}
 R_Q(N)=\sum_{k=1}^N\frac{Q(k)}{k^3},\qquad
 S_Q(N)=\sum_{k=1}^N\frac{Q(k)}{k^2}.
\end{equation}
They are polynomials because $Q$ is divisible by $X^4$.

\begin{theorem}[All odd-degree mixed coefficients]\label{thm:pair}
With the preceding definitions,
\begin{equation}\label{eq:pair-master}
 \boxed{\mathscr M(P)=p_2\zeta(3)+\lambda(P)\zeta(2)+\rho(Q),}
\end{equation}
where the rational number $\rho(Q)$ is given by the finite
coefficient extraction
\begin{equation}\label{eq:pair-rho}
 \rho(Q)=\CT_N\left[
 Q(N)A_NB_N+R_Q(N)A_N+S_Q(N)B_N\right].
\end{equation}
The asymptotic expansions of $A_N$ and $B_N$ at infinity make
\eqref{eq:pair-rho} a finite Bernoulli calculation.
In particular, $p_2=\lambda(P)=0$ gives a rational ordinary sum.
\end{theorem}

\begin{proof}
First assume $p_2=0$ and write $P=Q$. Define
$V_n=P(n)A_nB_n+R_P(n)A_n+S_P(n)B_n$.
The shifts $A_{n+1}=A_n-(n+1)^{-2}$ and
$B_{n+1}=B_n-(n+1)^{-3}$ give
\begin{equation}\label{eq:pair-V}
 V_{n+1}-V_n=W_P(x_n)A_nB_n+E_P(n+1),
\end{equation}
where
\[
 E_P(y)=\frac{P(y)}{y^5}-\frac{R_P(y)}{y^2}-\frac{S_P(y)}{y^3}.
\]
Faulhaber's formula shows that this Laurent polynomial has no
$y^{-1}$ term and that its $y^{-2}$ coefficient is $-\lambda(P)$.
For completeness, the linear coefficient of $R_P$ is $p_4/2$,
the quadratic coefficient of $S_P$ is $p_4/2$, and the linear
coefficient of $S_P$ is $\lambda(P)$. The possible $y^{-1}$
coefficient is therefore $p_4-p_4/2-p_4/2=0$.
Thus $E_P(y)=E_P^{\rm pol}(y)-\lambda(P)y^{-2}$.

Sum \eqref{eq:pair-V} for $0\le n<N$ and take cutoff constants.
The finite sum of the polynomial $E_P^{\rm pol}(k)$, $1\le k\le N$,
is a polynomial in $N$ with zero constant. Also $V_0=0$ and
$H_N^{(2)}\to\zeta(2)$. The ordinary-sum normalization from
Proposition~\ref{prop:regular} now gives
$\mathscr M(P)=\rho(P)+\lambda(P)\zeta(2)$.

For $P(X)=X^2$, finite summation by parts reduces the value to
\[
 \sum_{k\ge1}\frac{A_k}{k}+\sum_{k\ge1}B_k+\zeta(3)
 =\zeta(2,1)+\zeta(2)-\zeta(3)+\zeta(3)
 =\zeta(2)+\zeta(3).
\]
We used the classical Euler identity $\zeta(2,1)=\zeta(3)$.
Adding $p_2$ times this case proves \eqref{eq:pair-master}.
\end{proof}

Taking $P=F_1,F_3,F_5$ gives, respectively,
\begin{align}
 \sum_{n\ge0}x_nA_nB_n
   &={}\frac{\zeta(2)+\zeta(3)}2,\label{eq:pair1}\\
 \sum_{n\ge0}\left(x_n^3A_nB_n-\frac12\right)
   &={}-\frac{\zeta(2)}{12}-\frac{\zeta(3)}8-\frac5{48},\label{eq:pair3}\\
 \sum_{n\ge0}\left(x_n^5A_nB_n-\frac{x_n^2}{2}+\frac16\right)
   &={}\frac{47\zeta(2)}{1440}+\frac{7\zeta(3)}{96}+\frac{11}{288}.
 \label{eq:pair5}
\end{align}
Theorem~\ref{thm:pair} is an all-degree result; these examples are
not the basis of an extrapolation.

\subsection{A direct finite proof of the rational identity}

The combination $23F_1+232F_3+240F_5$ simplifies to
$P(X)=40X^6+8X^4$. Its $p_2$ is zero and
$\lambda(P)=40B_4+8B_2=0$. More explicitly, define
\begin{equation}\label{eq:telescoper}
\begin{split}
 U_n={}&8n^4(5n^2+1)A_nB_n
 +2n(n+1)(5n^2+5n+2)A_n\\
 &+4n^2(n+1)^2(2n+1)B_n-34n^3+9n^2-33n.
\end{split}
\end{equation}
Direct substitution of the two tail shifts gives the exact identity
\begin{equation}\label{eq:telescoper-difference}
 U_{n+1}-U_n=(23x_n+232x_n^3+240x_n^5)A_nB_n-120x_n^2-76.
\end{equation}
Here $U_0=0$. The ordinary uncentered expansions are
\begin{align*}
 A_n&=n^{-1}-\tfrac12n^{-2}+\tfrac16n^{-3}
                      -\tfrac1{30}n^{-5}+O(n^{-7}),\\
 B_n&=\tfrac12n^{-2}-\tfrac12n^{-3}+\tfrac14n^{-4}
                      -\tfrac1{12}n^{-6}+O(n^{-8}).
\end{align*}
They give $U_n=-15+11/(3n)+O(n^{-2})$. Summing
\eqref{eq:telescoper-difference} proves \eqref{eq:intro-rational}.
Absolute convergence also follows directly from the centered
subtraction, whose remainder is $O(x_n^{-2})$.

Equivalently, in polygamma notation,
\begin{equation}\label{eq:polygamma-rational}
\begin{split}
 \sum_{n\ge0}\bigl[&(23x_n+232x_n^3+240x_n^5)
                \psi'(n+1)\psi^{(2)}(n+1)\\
                &\hspace{25mm}+240x_n^2+152\bigr]=30.
\end{split}
\end{equation}
The sign here follows from $B_n=-\psi^{(2)}(n+1)/2$.

\subsection{Some further genuinely mixed examples}

The same finite compiler gives the following exact identities.
They illustrate odd-order participation and are accompanied by
full coefficient records in \sref{results/mixed\_moments.json}:
\begin{align}
 \sum_{n\ge0}A_n^2B_n
   &={}\frac32\zeta(3)^2+\frac76\zeta(6),\label{eq:mixed-cubic0}\\
 \sum_{n\ge0}x_n^2A_n^2B_n
   &={}\frac13\zeta(2)+\frac43\zeta(3)-\frac7{12}\zeta(4)
                   -\frac18\zeta(3)^2-\frac7{72}\zeta(6),\label{eq:mixed-cubic2}\\
 \sum_{n\ge0}\left(x_n^4A_n^2B_n-\frac12\right)
   &={}-\frac3{20}\zeta(2)-\frac7{15}\zeta(3)+\frac7{24}\zeta(4)
       +\frac7{160}\zeta(3)^2+\frac{49}{1440}\zeta(6)-\frac3{40}.
 \label{eq:mixed-cubic4}
\end{align}
With $C_n=\zeta(4,n+1)=\psi^{(3)}(n+1)/6$, an unweighted
weight-eight example is
\begin{equation}\label{eq:ABC}
 \sum_{n\ge0}A_nB_nC_n
 =-\frac{15}{2}\zeta(2)\zeta(3)^2+30\zeta(3)\zeta(5)
                   +\frac92\zeta(5,3)-\frac{1261}{72}\zeta(8).
\end{equation}
All unsubtracted sums displayed here converge absolutely.
The proofs are Theorem~\ref{thm:master} followed by the finite
classical relation system of Appendix~\ref{app:certificates};
the exact verifier regenerates these outputs. The unweighted
index-ordering principle itself is not new \cite{Hoffman}.

\section{A finite polylogarithm kernel and final-index zeta derivatives}
\label{sec:polylog}

Return to an arbitrary tuple $\boldsymbol p$. For $|z|<1$ set
\[
 \Li_{w_1,\ldots,w_\ell}(z)
 =\sum_{n_1>\cdots>n_\ell\ge1}
       \frac{z^{n_1}}{n_1^{w_1}\cdots n_\ell^{w_\ell}},
 \qquad \Li_{\varnothing}(z)=1,
\]
and use the coefficients from \eqref{eq:tail-harmonic-expansion}
to define the finite kernel
\begin{equation}\label{eq:Fkernel}
 \mathcal F_{\boldsymbol p}(z)=\sum_{\boldsymbol w}
                                  b_{\boldsymbol w}\Li_{\boldsymbol w}(z).
\end{equation}
Its depth is at most $d$ in this explicit representation; no minimal
functional depth is claimed.

\begin{proposition}[The ordinary generating function]\label{prop:kernel}
For $|z|<1$,
\begin{equation}\label{eq:tailGF}
 \boxed{\sum_{n\ge0}T_nz^n=\frac{\mathcal F_{\boldsymbol p}(z)}{1-z}.}
\end{equation}
For example,
\begin{equation}\label{eq:F23}
\begin{split}
 \mathcal F_{2,3}(z)={}&\zeta(2)\zeta(3)-\zeta(3)\Li_2(z)
 -\zeta(2)\Li_3(z)\\
 &+\Li_{2,3}(z)+\Li_{3,2}(z)+\Li_5(z).
\end{split}
\end{equation}
\end{proposition}
\begin{proof}
The absolutely convergent identity
$\sum_{n\ge0}H_n(\boldsymbol w)z^n=\Li_{\boldsymbol w}(z)/(1-z)$
follows by summing first over $n\ge n_1$. Apply it to the finite
expansion \eqref{eq:tail-harmonic-expansion}. Formula
\eqref{eq:F23} uses $H_n^{(2)}H_n^{(3)}=H_n(2,3)+H_n(3,2)+H_n^{(5)}$.
\end{proof}

Define a holomorphic Dirichlet family on $\Re s>1-R$ by
\begin{equation}\label{eq:Edef}
 \mathcal E_{\boldsymbol p}(s)=\sum_{n\ge1}\frac{T_{n-1}}{n^s}.
\end{equation}
For an ordered set partition $\pi=(B_1,\ldots,B_\ell)$ of
$\{1,\ldots,d\}$, let $w_{\pi,j}=\sum_{i\in B_j}p_i$.
Every ordered partition is included once, even when different
partitions give the same word.

\begin{theorem}[MZV and Mellin realizations]\label{thm:E}
On $\Re s>1-R$,
\begin{equation}\label{eq:E-MZV}
 \mathcal E_{\boldsymbol p}(s)=\sum_{\pi}\left[
 \zeta(w_{\pi,1},\ldots,w_{\pi,\ell},s)
 +\zeta(w_{\pi,1},\ldots,w_{\pi,\ell-1},w_{\pi,\ell}+s)
 \right].
\end{equation}
Every series on the right converges there. On $\Re s>0$ it also has
the ordinary integral representation
\begin{equation}\label{eq:E-integral}
 \boxed{\mathcal E_{\boldsymbol p}(s)=\frac1{\Gamma(s)}
 \int_0^1\frac{\mathcal F_{\boldsymbol p}(z)}{1-z}
                                  (-\log z)^{s-1}\dd z.}
\end{equation}
The identities may be differentiated any finite number of times
in $s$ on compact subsets of their respective domains.
\end{theorem}

\begin{proof}
In $T_{n-1}$ every tail index is at least $n$. Group equal tail
indices and order their blocks decreasingly. The distinguished
index $n$ is either strictly below the last block or equal to it.
These are exactly the two terms in \eqref{eq:E-MZV}. Absolute
convergence follows from $T_{n-1}=O(n^{-R})$; the same bounds with
fixed powers of $\log n$ give local uniformity of all derivatives.
The individual multiple series satisfy the same convergence
conditions: all proper initial block sums are at least twice
the number of blocks, and the full sum satisfies the condition
because $\Re s>1-R$.

For the integral, use
$n^{-s}=\Gamma(s)^{-1}\int_0^1z^{n-1}(-\log z)^{s-1}\dd z$
and \eqref{eq:tailGF}. Absolute convergence justifies the exchange.
Near $z=1$, the generating function is bounded when $R>1$ and
is at most logarithmic when $R=1$. Near zero it is bounded.
These endpoint bounds also justify all order derivatives for
$\Re s>0$.
\end{proof}

At $s=1$ a particularly useful all-index form uses the polynomials
\begin{equation}\label{eq:Gamma-Appell}
 \mathscr P_q(y)=q![u^q]\frac{e^{uy}}{\Gamma(1+u)}.
\end{equation}
For example,
\[
 \mathscr P_0=1,\quad \mathscr P_1=y+\gamma,\quad
 \mathscr P_2=(y+\gamma)^2-\zeta(2),
\]
\[
 \mathscr P_3=(y+\gamma)^3-3\zeta(2)(y+\gamma)+2\zeta(3).
\]
Differentiating \eqref{eq:E-integral} gives
\begin{equation}\label{eq:E-jets}
 \mathcal E_{\boldsymbol p}^{(q)}(1)=
 \int_0^1\frac{\mathcal F_{\boldsymbol p}(z)}{1-z}
                         \mathscr P_q(\log(-\log z))\dd z.
\end{equation}
All cancellations in $\mathcal F$ are kept inside the integral.
The reciprocal-Gamma coefficient construction is classical; its
use here fixes the constants in the tail-weighted Stieltjes identities.

\section{Stieltjes sums, polygamma derivatives, and anchored primitives}
\label{sec:transport}

Let
\begin{equation}\label{eq:Delta}
 \Delta_n=T_{n-1}-T_n\quad(n\ge1),\qquad
 Z_{\boldsymbol p}=T_0=\prod_i\zeta(p_i).
\end{equation}
These positive weights satisfy
\begin{equation}\label{eq:Delta-properties}
 \sum_{n\ge1}\Delta_n=Z_{\boldsymbol p},\qquad
 \sum_{n>k}\Delta_n=T_k,\qquad
 \Delta_n=O(n^{-R-1}).
\end{equation}
Only these elementary tail properties are required for the first
transport theorem. They follow either from telescoping or from
$\zeta(p,n)-\zeta(p,n+1)=n^{-p}$.

\begin{theorem}[Every Stieltjes index and argument]\label{thm:Stieltjes}
For every $q\ge0$ and $a>0$,
\begin{equation}\label{eq:Stieltjes}
 \boxed{\sum_{n\ge1}\Delta_n\gamma_q(n+a)
 =Z_{\boldsymbol p}\gamma_q(a)
       -\sum_{k\ge0}T_k\frac{\log^q(k+a)}{k+a}.}
\end{equation}
Both sums converge absolutely, locally uniformly for $a>0$.
At $a=1$ this becomes the convergent final-index derivative identity
\begin{equation}\label{eq:Stieltjes-MZV}
 \sum_{n\ge1}\Delta_n\gamma_q(n+1)
 =Z_{\boldsymbol p}\gamma_q(1)
                 -(-1)^q\mathcal E_{\boldsymbol p}^{(q)}(1).
\end{equation}
The finite MZV expression \eqref{eq:E-MZV} and the polylogarithm
integral \eqref{eq:E-jets} specify the derivative on the right.
\end{theorem}

\begin{proof}
Comparing Laurent coefficients in the Hurwitz shift equation gives
\begin{equation}\label{eq:gamma-shift}
 \gamma_q(n+a)=\gamma_q(a)
           -\sum_{k=0}^{n-1}\frac{\log^q(k+a)}{k+a}.
\end{equation}
Multiply by $\Delta_n$, sum, and use \eqref{eq:Delta-properties}.
The double sum of absolute values is bounded by
$\sum_k T_k|\log^q(k+a)|/(k+a)$, which converges because $R\ge1$.
Also $\gamma_q(n+a)=O(1+\log^{q+1}(n+1))$, either directly from
\eqref{eq:gamma-shift} or from Hurwitz asymptotics. Hence the left
sum is absolutely convergent. The same estimates are uniform
for $a$ in any compact subinterval of $(0,\infty)$.
Finally, differentiation of \eqref{eq:Edef} at $s=1$ gives
\eqref{eq:Stieltjes-MZV}.
\end{proof}

For an elementary check, take one tail of order two, so
$\Delta_n=n^{-2}$. The $q=0$ identity reduces to the classical
Euler sum $\sum H_n/n^2=2\zeta(3)$. At $q=1$ it gives the
ordinary derivative identity
\begin{equation}\label{eq:gamma1-example}
 \sum_{n\ge1}\frac{\gamma_1(n+1)}{n^2}
 =\zeta(2)\gamma_1(1)
   +\left.\frac{\partial}{\partial s}\zeta(2,s)\right|_{s=1}
   +\zeta'(3).
\end{equation}
No singular multiple-zeta regularization is needed: the differentiated
series converges in a neighborhood of $s=1$.

\begin{corollary}[Every polygamma order]\label{cor:polygamma}
For $r\ge0$ and $a>0$,
\begin{equation}\label{eq:polygamma-transport}
 \sum_{n\ge1}\Delta_n\psi^{(r)}(n+a)
 =Z_{\boldsymbol p}\psi^{(r)}(a)
       +(-1)^rr!\sum_{k\ge0}\frac{T_k}{(k+a)^{r+1}}.
\end{equation}
\end{corollary}
\begin{proof}
For $r=0$ this is \eqref{eq:Stieltjes} with $q=0$ and
$\gamma_0=-\psi$. For general $r$, use the differentiated digamma
shift equation and the same absolutely convergent interchange.
\end{proof}

\subsection{Gamma ratios and all Stieltjes primitives}

\begin{theorem}[Anchored Gamma and Hurwitz-zero primitives]\label{thm:primitives}
For $a,b>0$,
\begin{equation}\label{eq:Gamma-ratio}
\begin{split}
 \sum_{n\ge1}\Delta_n\log\frac{\Gamma(n+b)}{\Gamma(n+a)}
 ={}&Z_{\boldsymbol p}\log\frac{\Gamma(b)}{\Gamma(a)}
       +\sum_{k\ge0}T_k\log\frac{k+b}{k+a}.
\end{split}
\end{equation}
For $q\ge0$ set
\[
 \mathcal J_q(a,b)=\left.\partial_s^{q+1}\zeta(s,b)\right|_{s=0}
             -\left.\partial_s^{q+1}\zeta(s,a)\right|_{s=0}.
\]
Then
\begin{equation}\label{eq:zero-primitive}
\begin{split}
 &\sum_{n\ge1}\Delta_n\mathcal J_q(n+a,n+b)-Z_{\boldsymbol p}\mathcal J_q(a,b)\\
 &\hspace{10mm}=(-1)^q\sum_{k\ge0}T_k
          \left[\log^{q+1}(k+b)-\log^{q+1}(k+a)\right].
\end{split}
\end{equation}
All these series converge absolutely, and their anchors vanish when
$a=b$. The $q=0$ case of \eqref{eq:zero-primitive} is
\eqref{eq:Gamma-ratio}.
\end{theorem}

\begin{proof}
The finite Gamma product gives
\[
 \log\frac{\Gamma(n+b)}{\Gamma(n+a)}
 =\log\frac{\Gamma(b)}{\Gamma(a)}
                   +\sum_{k=0}^{n-1}\log\frac{k+b}{k+a}.
\]
The inner difference is $O((k+1)^{-1})$, so the same summation
argument as before proves \eqref{eq:Gamma-ratio}.

To prove the all-index primitive identity, differentiate
$\partial_a\zeta(s,a)=-s\zeta(s+1,a)$ at $s=0$, using
\eqref{eq:stieltjes-convention}. One obtains
\begin{equation}\label{eq:zero-derivative}
 \partial_a\left.\partial_s^{q+1}\zeta(s,a)\right|_{s=0}
 =-(q+1)(-1)^q\gamma_q(a).
\end{equation}
Integrate \eqref{eq:Stieltjes} from $a$ to $b$ and multiply by
$-(q+1)(-1)^q$. The elementary logarithmic primitive gives
\eqref{eq:zero-primitive}. Uniform convergence on the compact
interval between $a$ and $b$ justifies the integration. Finally,
$\zeta'(0,a)=\log\Gamma(a)-\frac12\log(2\pi)$ gives the
$q=0$ identification \cite{DLMF}.
\end{proof}

\subsection{A local analytic Gamma generator with explicit coefficients}

Setting $a=1$, $b=1+z$ in \eqref{eq:Gamma-ratio} gives, initially
for real $|z|<1$ and then holomorphically in that disk,
\begin{equation}\label{eq:Gamma-generator}
\begin{split}
 &\sum_{n\ge1}\Delta_n\log\frac{\Gamma(n+1+z)}{\Gamma(n+1)}
                  -Z_{\boldsymbol p}\log\Gamma(1+z)\\
 &\hspace{15mm}=\sum_{m\ge1}\frac{(-1)^{m+1}}m
                                     \mathcal E_{\boldsymbol p}(m)z^m.
\end{split}
\end{equation}
Here the branch is the analytic logarithm fixed to zero at $z=0$.
The coefficients are the finite convergent MZV combinations
\eqref{eq:E-MZV}. Absolute convergence on $|z|\le r<1$ follows
from the logarithm series and $T_k=O((k+1)^{-R})$.
This gives an anchored antiderivative family rather than a formula
with unspecified additive constants.

\section{Proof audit, computational evidence, and integration}
\label{sec:audit}

\subsection{What the finite checks establish}

The delivered code has two logically distinct parts. The exact
certificate verifier parses only restricted rational/MZV arithmetic,
expands stuffle products, independently forms binary shuffles by
interleaving positions, and reconstructs every cited relation.
It checks that each target difference equals the stated rational
combination exactly. A changed coefficient is deliberately rejected
for each of the three main anchor certificates.

The mixed-moment tests independently check finite centered summation
identities at rational cutoffs, the tail-product stuffle expansion,
Euler-constant cancellation, the refined weight ceilings, agreement
with the repeated-tail expansion, and the separate one-tail Faulhaber
formula for the $(2,3)$ family. The $-15$ telescoper is checked as
an identity of rational functions in $n,A,B$, not merely on sampled
special-function values. The exact ranks of the stated finite row
matrices are also recomputed. Their ranks concern rational relations
on formal word coordinates, not dimensions of numerical period spaces.

The final replay passes 369 exact assertions, including three deliberate
corruption controls. Complete records are in \sref{results/exact.json}.
The shortened certificates retain exactly the required nonzero relation
rows; every retained row is reconstructed from its defining words.
The numerical suite contains 40 comparisons at 90 working
decimal digits, including the three anchor sums, multiple truncations
of $\zeta(5,3)$ and fifteen mixed moments, finite Stieltjes and Gamma
transport, and a polylogarithm integral. Their largest recorded
residual is below $10^{-68}$ in the delivered run.

These numerical comparisons are diagnostics, \emph{not} interval
certificates. Tail acceleration uses finite asymptotic expansions
at two separate truncation choices. The analytic proofs supply
convergence and equality; decimal agreement supplies an independent
check of signs, constants, and normalization. There is no Lean or
other proof-assistant formalization in this package.

\subsection{A real normalization hazard, not an attributed source error}

The exclusive cutoff matters. One has
\[
 h_{(0)}(N)=N-1,\qquad \CT_N h_{(0)}(N)=-1,
\]
whereas the inclusive cutoff $\sum_{n=1}^{N}1=N$ has constant zero.
Replacing \eqref{eq:C-empty} by the inclusive value would silently
change the finite moment compiler. The separate check
$M_{(2),1}=-1/4$ detects such mistakes. This warning is not an
allegation that the inspected predecessor made that substitution.

No mathematical error in the specific predecessor formulas used
here was found in this audit. A confirmed editorial issue in the
inspected canonical README is the verbatim repetition of the paragraph
beginning with the complete 1,032-input first-letter-$X$ family and
its rank 1,008/kernel dimension 24. The proposed correction is simply
to remove the duplicate copy if it remains in the integration baseline.
The more important attribution correction in this research process
was to retain the newer report's ownership of the all-power recurrence
rather than relabel that theorem as a new result.

\subsection{Repository placement}

A natural report destination is the continuation
\sref{mixed-tail-anchor-reductions} under the existing
\sref{stieltjes-harmonic-resonance} spine, with a cross-reference
from the mixed-tail material of the lattice-contact report.
The package contains a short integration fragment, a claim-status
map, and source-audit notes. It does not modify the remote repository,
apply a canonical-manuscript patch, or change the status of the
Gaussian $S_6$, revised $S_8$, or general unequal-tail reconstruction
problems. Canonical integration should preserve the attribution
boundaries set out in Section~\ref{sec:scope}.

\section{Further research questions}
\label{sec:questions}

\paragraph{1. Reduce the remaining weight-eleven sextic anchor.}
The explicit value $R_{11}=\mathcal Z_{6,1}$ remains an admissible
finite MZV combination. A natural next task is to reduce it into
chosen weight-eleven coordinates, with a compact exact certificate,
and then rewrite the entire sextic generator in those coordinates.
No numerical independence is required to perform this reduction.

\paragraph{2. Find a mixed-order analytic generator with all anchors fixed.}
The compiler proves a finite arithmetic span for arbitrary integer
orders, including odd orders. The even-order source has stronger
analytic generating-function recurrences. Construct an equally
explicit recurrence for arbitrary mixed parity, determine its
necessary boundary terms, and give its finite multiple-polylogarithm
realization with all endpoint constants specified.

\paragraph{3. Classify rational polynomial combinations.}
For the $(2,3)$ family the two conditions $p_2=\lambda(P)=0$
produce rational values, with a finite telescoping construction.
Find analogous coefficient conditions for other mixed products.
The first concrete candidates are $(2,2,3)$ and $(2,3,4)$, whose
small formal coordinate spaces are visible in the supplied tables.
This is a rational coefficient-space problem, not a problem of
proving arithmetic independence.

\paragraph{4. Introduce rational and unequal argument shifts.}
Replace $\zeta(p_i,n+1)$ by $\zeta(p_i,n+a_i)$, first for positive
rational $a_i$. The centered polynomial subtraction no longer has
all its Hurwitz corrections equal to zero. Determine a finite
colored-MZV or Lerch-coordinate compiler retaining every correction,
and identify which parity-refined weight statements survive.

\paragraph{5. Differentiate the tail orders themselves.}
The Stieltjes transport already gives final-index derivatives of
MZVs. Differentiating the individual orders $p_i$ would introduce
logarithmic decorations on all ordered blocks. Prove a joint
continuation and differentiation theorem with explicit correction
terms at any new resonant parameter values. Finite experimental
agreement should not substitute for that local-uniform analysis.

\paragraph{6. Formalize the finite and analytic layers separately.}
The restricted rational certificate checker and the finite summation
identity \eqref{eq:finite-centered} are suitable first targets for
formalization. The next layer is the asymptotic remainder theorem
and its conversion of cutoff constants into ordinary subtracted
sums. Keeping these layers separate would prevent a finite rank
certificate from being mistaken for an analytic or transcendence proof.

\appendix
\section{Complete finite certificates for the three anchors}
\label{app:certificates}

This appendix makes the finite proof of Theorem~\ref{thm:anchors}
reconstructible without any numerical MZV database. All row indices
below are one-based. For weight $w$, list the admissible compositions
of $w$ in the order generated by choosing the first part increasingly,
then the next part increasingly, recursively. This gives $2^{w-2}$
columns, one for each formal MZV word.

For an index word $u=(u_1,\ldots,u_r)$, its binary encoding is
$0^{u_1-1}1\cdots0^{u_r-1}1$. Decode any binary word ending in $1$
by counting the zero runs plus one. Let $u*v$ be the positive
stuffle product of index words, and let $u\sha v$ denote the sum
obtained by shuffling their binary encodings and decoding. Here
$\sha$ is written in the tables' generating rule as the ordinary
shuffle operation; repeated words retain their multiplicities.

The ordered relation list is constructed as follows. For
$k=1,\ldots,w-2$, take $u=(1)$ if $k=1$, and otherwise take every
admissible composition $u$ of $k$ in the preceding order. For each
$u$, take every admissible composition $v$ of $w-k$ in that order.
Append the coefficient vector of
\begin{equation}\label{eq:DSrow}
 \mathcal R(u,v)=\zeta(u*v)-\zeta(u\mathbin{\shuffle}v),
\end{equation}
omitting a row only when its entire coefficient vector is zero.
After these rows, append $\zeta(u)-\zeta(u^\vee)$ for every
admissible word $u$ of weight $w$, in the same order, omitting
self-dual words. The dual $u^\vee$ is obtained by reversing the
binary encoding and exchanging zero and one. Duplicate nonzero
rows are not removed.

For convergent $u,v$, \eqref{eq:DSrow} is the equality between
the series product and the iterated-integral product. For $u=(1)$,
the same row is the one-leading-one regularized double-shuffle
identity: the divergent leading word cancels and the surviving
combination is admissible. This special case can be obtained by
subtracting the common logarithmic term before taking the endpoint
limit; it is also a direct instance of \cite{IKZ}. Duality follows
from reversal and $t\mapsto1-t$ in the iterated integrals.
Thus every row used is a proved identity.

\begin{center}
\begin{tabular}{rrrrr}
\toprule
Weight & Columns & Rows & Rational row rank & Rows in anchor certificate\\
\midrule
7&32&80&29&3\\
8&64&168&60&17\\
9&128&384&123&54\\
\bottomrule
\end{tabular}
\end{center}

Let $D_7,D_8,D_9$ be the left side minus the right side of
\eqref{eq:J7}, \eqref{eq:K8}, \eqref{eq:K9}, respectively,
after the anchor definitions and all stuffle products are expanded.
The tables give the exact vector identities
\begin{equation}\label{eq:certidentity}
 D_w=\sum_j c_j\mathcal R_{w,j}.
\end{equation}
All unlisted coefficients are zero. Ranks in the preceding table
are only diagnostics of this finite relation system; equality in
\eqref{eq:certidentity}, not a completeness conjecture, proves
the three anchor identities.

\subsection*{Weight 7 certificate}
\begin{center}
\small
\begin{tabular}{rr@{\qquad}rr@{\qquad}rr}
\toprule
$j$&$c_j$&$j$&$c_j$&$j$&$c_j$\\
\midrule

16&$-10$&24&$3$&32&$-1$\\[3pt]

\bottomrule
\end{tabular}
\end{center}

\subsection*{Weight 8 certificate}
\begin{center}
\small
\begin{tabular}{rr@{\qquad}rr@{\qquad}rr}
\toprule
$j$&$c_j$&$j$&$c_j$&$j$&$c_j$\\
\midrule

24&$\frac{85}{33}$&28&$\frac{85}{11}$&30&$\frac{265}{33}$\\[3pt]

31&$- \frac{125}{11}$&32&$\frac{635}{132}$&40&$\frac{415}{132}$\\[3pt]

44&$\frac{995}{132}$&46&$\frac{301}{132}$&47&$- \frac{68}{33}$\\[3pt]

48&$- \frac{1347}{88}$&56&$- \frac{250}{33}$&60&$- \frac{665}{132}$\\[3pt]

62&$- \frac{71}{12}$&63&$\frac{343}{66}$&64&$\frac{5239}{264}$\\[3pt]

72&$\frac{895}{132}$&80&$- \frac{951}{88}$&&\\[3pt]

\bottomrule
\end{tabular}
\end{center}

\subsection*{Weight 9 certificate}
\begin{center}
\small
\begin{tabular}{rr@{\qquad}rr@{\qquad}rr}
\toprule
$j$&$c_j$&$j$&$c_j$&$j$&$c_j$\\
\midrule

30&$-480$&31&$-960$&32&$80$\\[3pt]

44&$40$&46&$- \frac{985544}{417}$&47&$- \frac{1276904}{139}$\\[3pt]

48&$\frac{17615752}{5421}$&52&$120$&54&$- \frac{560744}{139}$\\[3pt]

55&$- \frac{1753280}{139}$&56&$\frac{9541218}{1807}$&58&$- \frac{696688}{139}$\\[3pt]

59&$- \frac{1545888}{139}$&60&$\frac{6447172}{1807}$&61&$- \frac{2298480}{139}$\\[3pt]

62&$- \frac{31629680}{5421}$&63&$- \frac{18416956}{1807}$&64&$- \frac{12669404}{5421}$\\[3pt]

76&$-50$&78&$\frac{224942}{417}$&79&$\frac{185462}{139}$\\[3pt]

80&$\frac{743996}{5421}$&84&$- \frac{280}{3}$&86&$\frac{660508}{417}$\\[3pt]

87&$\frac{346688}{139}$&88&$- \frac{434087}{2224}$&90&$\frac{72448}{139}$\\[3pt]

91&$- \frac{322884}{139}$&92&$- \frac{9854045}{7228}$&93&$- \frac{1463144}{139}$\\[3pt]

94&$\frac{7674219}{14456}$&95&$\frac{2933377}{834}$&96&$\frac{259435535}{86736}$\\[3pt]

111&$\frac{258000}{139}$&112&$- \frac{3555872}{1807}$&116&$80$\\[3pt]

118&$\frac{563662}{417}$&119&$- \frac{2868}{139}$&120&$- \frac{3082043}{6672}$\\[3pt]

122&$- \frac{18272}{417}$&123&$- \frac{507776}{417}$&124&$- \frac{22787469}{14456}$\\[3pt]

125&$- \frac{723812}{139}$&126&$- \frac{6757760}{5421}$&127&$- \frac{6947744}{5421}$\\[3pt]

128&$- \frac{29083111}{86736}$&142&$- \frac{402598}{417}$&144&$- \frac{3716515}{1807}$\\[3pt]

150&$\frac{157344}{139}$&151&$\frac{879564}{139}$&156&$\frac{185281999}{86736}$\\[3pt]

158&$\frac{166192045}{86736}$&159&$\frac{2900207}{7228}$&160&$- \frac{5417647}{10842}$\\[3pt]

\bottomrule
\end{tabular}
\end{center}


\section{Reproduction and artifact contents}
\label{app:repro}

The standalone source is \sref{article.tex}; the PDF is compiled
from it without external figures or a bibliography database.
The Python dependencies used in the delivered run are recorded
in \sref{requirements.txt} and \sref{results/environment.json}.
From the package directory, run
\begin{verbatim}
python verification/run_all.py
python build.py
\end{verbatim}
The exact verifier does not need network access or a repository
checkout. The numerical comparisons use mpmath. The complete
coefficient certificates are in \sref{certificates/}; the independent
verifier accepts only restricted arithmetic syntax when reading them.
The generated example table is \sref{results/mixed\_moments.json}.
Original delivery evidence is in \sref{results/}; rerunning a suite
rewrites its corresponding result file, not repository source files.

\begin{thebibliography}{9}
\bibitem{DLMF}
F. W. J. Olver et al., eds., \emph{NIST Digital Library of Mathematical
Functions}, Sections 25.11 and 25.12, Hurwitz zeta and polylogarithms.
\href{https://dlmf.nist.gov/25.11}{dlmf.nist.gov/25.11} and
\href{https://dlmf.nist.gov/25.12}{dlmf.nist.gov/25.12}.

\bibitem{Hoffman}
M. E. Hoffman, Sums of products of Riemann zeta tails,
\emph{Mediterranean Journal of Mathematics} \textbf{13} (2016),
2771--2781. DOI: 10.1007/s00009-015-0653-9.
\href{https://arxiv.org/abs/1602.03399}{arXiv:1602.03399}.

\bibitem{IKZ}
K. Ihara, M. Kaneko, and D. Zagier,
Derivation and double shuffle relations for multiple zeta values,
\emph{Compositio Mathematica} \textbf{142} (2006), 307--338.
\href{https://doi.org/10.1112/S0010437X0500182X}{DOI:10.1112/S0010437X0500182X}.

\bibitem{Repo}
ProveIt Contributors, \emph{Polylogarithms and their Arithmetic Bridges},
canonical manuscript and status README,
\href{https://github.com/VladimirReshetnikov/ProveIt/tree/main/Analysis/Polylogarithms/docs/manuscript}{ProveIt polylogarithm manuscript}.
The exact scope of this article's inspection is recorded in
\sref{SOURCE\_AUDIT.md}.

\bibitem{Quartic}
ProveIt research continuation, \emph{Harmonic-Tail Reconstruction,
Quartic Trigamma Moments, and Spectral Stieltjes Calculus},
11 October 2026. Report continuation
\sref{quartic-spectral-reconstruction} in the
\sref{stieltjes-harmonic-resonance} spine.
\href{https://github.com/VladimirReshetnikov/ProveIt/tree/0dfef47533f8bbf23d12aa4f1894b9bea61f6e0a/Analysis/Polylogarithms/docs/reports/stieltjes-harmonic-resonance/continuations/quartic-spectral-reconstruction}{Pinned repository location}.

\bibitem{Bernoulli}
ProveIt research continuation, \emph{Bernoulli Cancellation,
Harmonic Gaps, and Spectral Faces}, 11 October 2026.
In particular, the all-power theorem and the fifth/sixth-power
anchor definitions and generators.
\href{https://github.com/VladimirReshetnikov/ProveIt/tree/0dfef47533f8bbf23d12aa4f1894b9bea61f6e0a/Analysis/Polylogarithms/docs/reports/stieltjes-harmonic-resonance/continuations/bernoulli-tails-spectral-faces}{Pinned repository location}.

\bibitem{Lattice}
ProveIt research continuation, \emph{Lattice Contacts and Mixed Zeta
Identities}, 11 October 2026. Unequal even-order tail recurrence
and its stated mixed generators.
\href{https://github.com/VladimirReshetnikov/ProveIt/tree/cc569dca8622dfc8fa02906da37d0a1b8831b5e7/Analysis/Polylogarithms/docs/reports/stieltjes-correlation-calculus/continuations/lattice-contacts-mixed-zeta}{Supplementary pinned repository location}.
\end{thebibliography}
\end{document}
